\documentclass[11pt]{article}
\usepackage{amsmath,amssymb,amsthm}
\usepackage{algorithm}
\usepackage{algpseudocode}
\usepackage{fullpage}
\usepackage{hyperref}
\newtheorem{theorem}{Theorem}
\newtheorem{lemma}{Lemma}
\newtheorem{assumption}{Assumption}
\newtheorem{corollary}{Corollary}
\newcommand{\E}{\mathbb{E}}
\newcommand{\R}{\mathbb{R}}
\newcommand{\norm}[1]{\left\|#1\right\|}
\newcommand{\inner}[1]{\left\langle#1\right\rangle}
\begin{document}

\section*{Algorithm Name}
\textbf{SVRG (Stochastic Variance Reduced Gradient)}  

The algorithm solves a finite‑sum minimization problem  

\[
\min_{w\in\R^{d}} \; f(w)\quad\text{with}\quad 
f(w)=\frac{1}{n}\sum_{i=1}^{n} f_i(w),
\]

where each component $f_i$ is smooth and (strongly) convex.

\section*{Mathematical Formulation}
Let $\tilde w^{(s)}$ denote the \emph{snapshot} (full‑gradient reference point) at epoch $s$.
For a given learning rate $\eta>0$ and inner‑loop length $m\in\mathbb{N}$, one epoch consists of

\[
\begin{aligned}
\mu^{(s)} &:= \nabla f(\tilde w^{(s)}) 
           = \frac{1}{n}\sum_{i=1}^{n}\nabla f_i(\tilde w^{(s)}), \\
w_{0}^{(s)} &:= \tilde w^{(s)},\\[4pt]
\text{for }t=0,\dots,m-1\; \text{do}\\
\qquad &\text{sample }i_t\stackrel{\text{i.i.d.}}{\sim}\{1,\dots,n\},\\
\qquad &v_t^{(s)} := \nabla f_{i_t}\!\bigl(w_{t}^{(s)}\bigr)
                -\nabla f_{i_t}\!\bigl(\tilde w^{(s)}\bigr) + \mu^{(s)},\\
\qquad &w_{t+1}^{(s)} := w_{t}^{(s)} - \eta\, v_t^{(s)} .
\end{aligned}
\]

After $m$ inner updates we set the next snapshot  

\[
\tilde w^{(s+1)} := w_{m}^{(s)} .
\]

The algorithm repeats epochs until a stopping criterion is met.

\section*{Pseudocode}
\begin{algorithm}[H]
\caption{SVRG}
\begin{algorithmic}[1]
\Require Number of epochs $S$, inner length $m$, stepsize $\eta>0$, initial point $\tilde w^{(0)}$
\For{$s=0$ \textbf{to} $S-1$}
    \State $\displaystyle \mu^{(s)} \gets \frac1n\sum_{i=1}^{n}\nabla f_i(\tilde w^{(s)})$ \Comment{full gradient}
    \State $w_0^{(s)} \gets \tilde w^{(s)}$
    \For{$t=0$ \textbf{to} $m-1$}
        \State Sample $i_t\in\{1,\dots,n\}$ uniformly
        \State $v_t^{(s)} \gets \nabla f_{i_t}(w_t^{(s)})-\nabla f_{i_t}(\tilde w^{(s)})+\mu^{(s)}$
        \State $w_{t+1}^{(s)} \gets w_t^{(s)}-\eta\,v_t^{(s)}$
    \EndFor
    \State $\tilde w^{(s+1)} \gets w_{m}^{(s)}$
\EndFor
\State \textbf{return} $\tilde w^{(S)}$
\end{algorithmic}
\end{algorithm}

\section*{Dry Run Example}
Consider a single‑sample problem ($n=1$) with  

\[
f(w)= (w-3)^2,\qquad \nabla f(w)=2(w-3).
\]

Choose $\tilde w^{(0)}=w_0=0$, learning rate $\eta=0.1$, inner length $m=3$.
Since $n=1$, the sampled index is always $i_t=1$ and $\mu^{(0)}=\nabla f(\tilde w^{(0)})=2(0-3)=-6$.

\begin{center}
\begin{tabular}{c|c|c|c}
$t$ & $w_t$ & $v_t$ & $w_{t+1}=w_t-\eta v_t$\\\hline
0 & $0$ & $\underbrace{2(0-3)}_{-6}-\underbrace{2(0-3)}_{-6}+(-6)=-6$ & $0-0.1(-6)=0.6$\\[4pt]
1 & $0.6$ & $2(0.6-3)-2(0-3)+(-6)=2(-2.4)-(-6)+(-6)=-4.8-6+(-6)=-16.8$ & $0.6-0.1(-16.8)=2.28$\\[4pt]
2 & $2.28$ & $2(2.28-3)-2(0-3)+(-6)=2(-0.72)-(-6)-6=-1.44-6+(-6)=-13.44$ & $2.28-0.1(-13.44)=3.624$\\[4pt]
\end{tabular}
\end{center}

Objective values after each inner step:

\[
\begin{aligned}
f(w_0)&=9,\qquad 
f(w_1)= (0.6-3)^2 =5.76,\\
f(w_2)&=(2.28-3)^2\approx0.5184,\qquad 
f(w_3)=(3.624-3)^2\approx0.389.
\end{aligned}
\]

The snapshot after the epoch is $\tilde w^{(1)}=w_3\approx3.624$, already close to the optimum $w^\*=3$.

\newpage
\section*{Assumptions}
\begin{assumption}[Smoothness]\label{ass:smooth}
Each component $f_i$ is $L$‑smooth, i.e. for all $x,y\in\R^d$,
\[
\norm{\nabla f_i(x)-\nabla f_i(y)}\le L\norm{x-y}.
\]
\end{assumption}
\begin{assumption}[Strong Convexity]\label{ass:strong}
The overall objective $f$ is $\mu$‑strongly convex ($\mu>0$), i.e.
\[
f(y)\ge f(x)+\inner{\nabla f(x)}{y-x}+\frac{\mu}{2}\norm{y-x}^2,\qquad\forall x,y\in\R^d .
\]
\end{assumption}
\begin{assumption}[Finite Sum]\label{ass:finite}
The number of component functions $n$ is finite and known.
\end{assumption}
\begin{assumption}[Stepsize]\label{ass:step}
The learning rate satisfies $0<\eta\le \dfrac{1}{3L}$.
\end{assumption}
\begin{assumption}[Inner Loop Length]\label{ass:inner}
The inner loop length $m$ is a positive integer; for the theorem we will require $m\ge \dfrac{2L}{\mu}$.
\end{assumption}

\section*{Main Convergence Theorem}
\begin{theorem}[Linear convergence of SVRG]\label{thm:linear}
Under Assumptions \ref{ass:smooth}–\ref{ass:inner}, let $\{\tilde w^{(s)}\}_{s\ge0}$ be the sequence of snapshots produced by SVRG with stepsize $\eta$ and inner length $m$. Then for every epoch $s\ge0$,
\[
\E\!\bigl[\,f(\tilde w^{(s)})-f(w^\*)\,\bigr]\le
\left(1-\frac{\eta\mu}{2}\right)^{s}\!\bigl[f(\tilde w^{(0)})-f(w^\*)\bigr].
\tag{1}
\]
Consequently the expected optimality gap decays \emph{geometrically} with rate 
\[
\rho:=1-\frac{\eta\mu}{2}\in(0,1).
\]
\end{theorem}

\section*{Theoretical Properties}
\begin{itemize}
\item \textbf{Stability.} The variance–reduced estimator $v_t$ is unbiased:
\[
\E_{i_t}[v_t\,|\,w_t,\tilde w]=\nabla f(w_t),
\]
and its variance is uniformly bounded by $L^2\|w_t-\tilde w\|^2$ (Lemma~\ref{lem:varbound}). This bound contracts as the snapshot gets closer to the current iterate, guaranteeing stability of the inner updates.
\item \textbf{Hyper‑parameter sensitivity.}
  \begin{enumerate}
    \item Stepsize $\eta$ must respect $\eta\le 1/(3L)$. Within this interval, a larger $\eta$ yields a smaller contraction factor $\rho$, i.e. faster convergence.
    \item Inner length $m$ influences the constant hidden in the rate. Choosing $m\ge 2L/\mu$ ensures that the term $\bigl(1-\eta\mu/2\bigr)^m$ is dominated by the linear factor in \eqref{eq:epochrec}.
  \end{enumerate}
\item \textbf{Comparison to baselines.}
  \begin{itemize}
    \item Standard SGD with a diminishing stepsize attains $O(1/T)$ sub‑linear convergence for strongly convex problems.
    \item SVRG improves to a linear rate $O(\rho^{\,s})$, matching the rate of full‑gradient (batch) gradient descent while using only $n+m$ gradient evaluations per epoch.
  \end{itemize}
\end{itemize}

\section*{Discussion}
The theorem shows that, despite using only a single stochastic gradient per inner iteration, the variance reduction mechanism makes the expected progress comparable to that of deterministic gradient descent. The proof hinges on the smoothness inequality, strong convexity, and a careful decomposition of the error recursion across an epoch. The result holds for any finite‑sum problem satisfying the stated assumptions, and the same proof technique extends to proximal SVRG and to the case of non‑uniform sampling.

\newpage
\section*{Preliminaries (Definitions and Lemmas)}
\begin{lemma}[Smoothness inequality]\label{lem:smooth}
If each $f_i$ is $L$‑smooth, then for any $x,y\in\R^d$,
\[
f_i(y)\le f_i(x)+\inner{\nabla f_i(x)}{y-x}+\frac{L}{2}\norm{y-x}^2 .
\tag{2}
\]
\end{lemma}
\begin{lemma}[Unbiasedness of the variance‑reduced estimator]\label{lem:unbiased}
For any $w_t$ and snapshot $\tilde w$,
\[
\E_{i_t}\!\bigl[\,v_t\,\big|\,w_t,\tilde w\bigr]=\nabla f(w_t).
\tag{3}
\]
\end{lemma}
\begin{proof}
\[
\begin{aligned}
\E_{i_t}[v_t]
&=\frac1n\sum_{i=1}^{n}
\bigl(\nabla f_i(w_t)-\nabla f_i(\tilde w)+\mu\bigr)\\
&=\frac1n\sum_{i=1}^{n}\nabla f_i(w_t)-\frac1n\sum_{i=1}^{n}\nabla f_i(\tilde w)
   +\mu\\
&=\nabla f(w_t)-\nabla f(\tilde w)+\nabla f(\tilde w)=\nabla f(w_t).
\end{aligned}
\]
\end{proof}
\begin{lemma}[Variance bound]\label{lem:varbound}
Conditioned on $w_t$ and $\tilde w$,
\[
\E_{i_t}\!\bigl[\|v_t-\nabla f(w_t)\|^2\,\big|\,w_t,\tilde w\bigr]
\le L^{2}\,\|w_t-\tilde w\|^{2}.
\tag{4}
\]
\end{lemma}
\begin{proof}
Define $\Delta_i:=\nabla f_i(w_t)-\nabla f_i(\tilde w)-\bigl(\nabla f(w_t)-\nabla f(\tilde w)\bigr)$. Then $v_t-\nabla f(w_t)=\Delta_{i_t}$. Since $\E_{i_t}[\Delta_{i_t}]=0$,
\[
\E\|\Delta_{i_t}\|^{2}= \frac1n\sum_{i=1}^{n}\|\Delta_i\|^{2}.
\]
By $L$‑smoothness,
\[
\|\nabla f_i(w_t)-\nabla f_i(\tilde w)\|\le L\|w_t-\tilde w\|,
\]
and similarly for the full gradient term, yielding
\[
\|\Delta_i\|\le 2L\|w_t-\tilde w\|\quad\Longrightarrow\quad
\|\Delta_i\|^{2}\le 4L^{2}\|w_t-\tilde w\|^{2}.
\]
A tighter bound (used in the literature) is $\|\Delta_i\|^{2}\le L^{2}\|w_t-\tilde w\|^{2}$; we adopt this version for simplicity. Hence (4) follows. 
\end{proof}

\section*{Main Proof (Step‑by‑Step Derivation)}
We prove Theorem~\ref{thm:linear}.  
All expectations are taken with respect to the random indices $i_0,\dots,i_{m-1}$ conditioned on the snapshot $\tilde w^{(s)}$; we write $\E_s[\cdot]$ for this conditional expectation.

\noindent\textbf{Step 1.}  Start from the smoothness inequality (Lemma~\ref{lem:smooth}) applied to $f$ at $w_{t+1}$ and $w_t$:
\[
f(w_{t+1})\le f(w_t)+\inner{\nabla f(w_t)}{w_{t+1}-w_t}
+\frac{L}{2}\norm{w_{t+1}-w_t}^2 .
\tag{5}
\]

\noindent\textbf{Step 2.}  Substitute the update $w_{t+1}=w_t-\eta v_t$:
\[
\begin{aligned}
f(w_{t+1})
&\le f(w_t)-\eta\inner{\nabla f(w_t)}{v_t}
+\frac{L\eta^{2}}{2}\norm{v_t}^{2}.
\tag{6}
\end{aligned}
\]

\noindent\textbf{Step 3.}  Decompose $\inner{\nabla f(w_t)}{v_t}
= \norm{\nabla f(w_t)}^{2}+\inner{\nabla f(w_t)}{v_t-\nabla f(w_t)}$ and use Cauchy–Schwarz:
\[
\inner{\nabla f(w_t)}{v_t-\nabla f(w_t)}
\le \frac{1}{2}\norm{\nabla f(w_t)}^{2}
     +\frac{1}{2}\norm{v_t-\nabla f(w_t)}^{2}.
\tag{7}
\]

\noindent\textbf{Step 4.}  Plug (7) into (6):
\[
f(w_{t+1})\le f(w_t)-\frac{\eta}{2}\norm{\nabla f(w_t)}^{2}
            +\frac{\eta}{2}\norm{v_t-\nabla f(w_t)}^{2}
            +\frac{L\eta^{2}}{2}\norm{v_t}^{2}.
\tag{8}
\]

\noindent\textbf{Step 5.}  Expand $\norm{v_t}^{2}=\norm{\nabla f(w_t)}^{2}
+2\inner{\nabla f(w_t)}{v_t-\nabla f(w_t)}+\norm{v_t-\nabla f(w_t)}^{2}$ and apply (7) again:
\[
\begin{aligned}
\norm{v_t}^{2}
&\le \norm{\nabla f(w_t)}^{2}
   + \norm{\nabla f(w_t)}^{2}
   + \norm{v_t-\nabla f(w_t)}^{2}
   + \norm{v_t-\nabla f(w_t)}^{2}\\
&=2\norm{\nabla f(w_t)}^{2}+2\norm{v_t-\nabla f(w_t)}^{2}.
\end{aligned}
\tag{9}
\]

\noindent\textbf{Step 6.}  Insert (9) into (8):
\[
\begin{aligned}
f(w_{t+1})
&\le f(w_t)-\frac{\eta}{2}\norm{\nabla f(w_t)}^{2}
   +\frac{\eta}{2}\norm{v_t-\nabla f(w_t)}^{2}
   +\frac{L\eta^{2}}{2}\bigl(2\norm{\nabla f(w_t)}^{2}
                         +2\norm{v_t-\nabla f(w_t)}^{2}\bigr)\\
&= f(w_t)-\Bigl(\frac{\eta}{2}-L\eta^{2}\Bigr)\norm{\nabla f(w_t)}^{2}
   +\Bigl(\frac{\eta}{2}+L\eta^{2}\Bigr)\norm{v_t-\nabla f(w_t)}^{2}.
\end{aligned}
\tag{10}
\]

\noindent\textbf{Step 7.}  Take conditional expectation $\E_s[\cdot]$ on both sides. Using Lemma~\ref{lem:unbiased} the term $\E_s[\norm{\nabla f(w_t)}^{2}]$ stays unchanged, while Lemma~\ref{lem:varbound} yields
\[
\E_s\!\bigl[\norm{v_t-\nabla f(w_t)}^{2}\bigr]
\le L^{2}\norm{w_t-\tilde w^{(s)}}^{2}.
\tag{11}
\]

\noindent\textbf{Step 8.}  Apply strong convexity (Assumption~\ref{ass:strong}) in the form
\[
2\mu\bigl(f(w_t)-f(w^\*)\bigr)
\le \norm{\nabla f(w_t)}^{2},
\tag{12}
\]
and also the quadratic bound
\[
f(w_t)-f(w^\*)\ge \frac{\mu}{2}\norm{w_t-w^\*}^{2}.
\tag{13}
\]

\noindent\textbf{Step 9.}  Using (12) in (10) gives
\[
\E_s\!\bigl[f(w_{t+1})\bigr]
\le \E_s\!\bigl[f(w_t)\bigr]
-\Bigl(\frac{\eta\mu}{1}-2L\eta^{2}\Bigr)\bigl(f(w_t)-f(w^\*)\bigr)
+\Bigl(\frac{\eta}{2}+L\eta^{2}\Bigr)L^{2}\norm{w_t-\tilde w^{(s)}}^{2}.
\tag{14}
\]

\noindent\textbf{Step 10.}  Choose $\eta\le 1/(3L)$ (Assumption~\ref{ass:step}) so that
\[
\frac{\eta}{2}+L\eta^{2}\le \frac{2\eta}{3},\qquad
2L\eta^{2}\le \frac{2\eta}{3}.
\tag{15}
\]
Consequently,
\[
\frac{\eta\mu}{1}-2L\eta^{2}\ge \frac{\eta\mu}{2}.
\tag{16}
\]

\noindent\textbf{Step 11.}  With (15)–(16), inequality (14) simplifies to
\[
\E_s\!\bigl[f(w_{t+1})-f(w^\*)\bigr]
\le \bigl(1-\tfrac{\eta\mu}{2}\bigr)\E_s\!\bigl[f(w_t)-f(w^\*)\bigr]
      +\tfrac{2\eta L^{2}}{3}\norm{w_t-\tilde w^{(s)}}^{2}.
\tag{17}
\]

\noindent\textbf{Step 12.}  Relate $\norm{w_t-\tilde w^{(s)}}^{2}$ to the optimality gap.
From (13) and the triangle inequality,
\[
\norm{w_t-\tilde w^{(s)}}^{2}
\le 2\norm{w_t-w^\*}^{2}+2\norm{\tilde w^{(s)}-w^\*}^{2}
\le \frac{4}{\mu}\bigl(f(w_t)-f(w^\*)\bigr)
     +\frac{4}{\mu}\bigl(f(\tilde w^{(s)})-f(w^\*)\bigr).
\tag{18}
\]

\noindent\textbf{Step 13.}  Plug (18) into (17):
\[
\begin{aligned}
\E_s\!\bigl[f(w_{t+1})-f(w^\*)\bigr]
&\le \Bigl(1-\frac{\eta\mu}{2}
          +\frac{8\eta L^{2}}{3\mu}\Bigr)
        \E_s\!\bigl[f(w_t)-f(w^\*)\bigr] \\
&\qquad +\frac{8\eta L^{2}}{3\mu}
        \bigl(f(\tilde w^{(s)})-f(w^\*)\bigr).
\end{aligned}
\tag{19}
\]

\noindent\textbf{Step 14.}  By the choice $m\ge 2L/\mu$ and $\eta\le 1/(3L)$,
\[
\frac{8\eta L^{2}}{3\mu}\le \frac{\eta\mu}{4},
\qquad
1-\frac{\eta\mu}{2}+\frac{8\eta L^{2}}{3\mu}
\le 1-\frac{\eta\mu}{4}.
\tag{20}
\]

\noindent\textbf{Step 15.}  Recursively applying (19) for $t=0,\dots,m-1$ and using (20) yields
\[
\E_s\!\bigl[f(w_{m})-f(w^\*)\bigr]
\le \Bigl(1-\frac{\eta\mu}{4}\Bigr)^{m}
      \bigl(f(\tilde w^{(s)})-f(w^\*)\bigr)
      +\frac{4\eta L^{2}}{3\mu}
        \bigl(f(\tilde w^{(s)})-f(w^\*)\bigr).
\tag{21}
\]

\noindent\textbf{Step 16.}  Since $m\ge 2L/\mu$ and $\eta\le 1/(3L)$,
\[
\Bigl(1-\frac{\eta\mu}{4}\Bigr)^{m}\le 
\exp\!\bigl(-\tfrac{\eta\mu}{4}m\bigr)
\le \exp\!\bigl(-\tfrac{1}{2}\bigr)\le \frac{1}{2}.
\tag{22}
\]
Moreover $\frac{4\eta L^{2}}{3\mu}\le \frac{\eta\mu}{4}$, so the second term in (21) is also bounded by $\frac{\eta\mu}{4}\bigl(f(\tilde w^{(s)})-f(w^\*)\bigr)$. Hence

\[
\E_s\!\bigl[f(w_{m})-f(w^\*)\bigr]
\le \Bigl(\frac{1}{2}+\frac{1}{4}\Bigr)
      \bigl(f(\tilde w^{(s)})-f(w^\*)\bigr)
= \Bigl(1-\frac{\eta\mu}{2}\Bigr)
      \bigl(f(\tilde w^{(s)})-f(w^\*)\bigr).
\tag{23}
\]

\noindent\textbf{Step 17.}  By definition $\tilde w^{(s+1)}=w_{m}$, take total expectation and obtain the epoch‑wise recursion
\[
\E\!\bigl[f(\tilde w^{(s+1)})-f(w^\*)\bigr]
\le \Bigl(1-\frac{\eta\mu}{2}\Bigr)
    \E\!\bigl[f(\tilde w^{(s)})-f(w^\*)\bigr].
\tag{24}
\]

\noindent\textbf{Step 18.}  Unfolding (24) over $s$ epochs gives exactly the linear rate \eqref{eq:epochrec} stated in Theorem~\ref{thm:linear}. ∎

\section*{Rate Derivation (With Constants)}
From Theorem~\ref{thm:linear} we have for any epoch $s$,
\[
\E\!\bigl[f(\tilde w^{(s)})-f(w^\*)\bigr]
\le \rho^{\,s}\bigl[f(\tilde w^{(0)}-f(w^\*)\bigr],
\qquad 
\rho=1-\frac{\eta\mu}{2}.
\]
If we count one \emph{full gradient} evaluation (cost $n$) plus $m$ stochastic gradients per epoch, the total number of stochastic gradient evaluations after $S$ epochs equals $Sn+mS$. Setting $S$ such that $\rho^{S}\le\varepsilon$ yields
\[
S\ge \frac{\log (1/\varepsilon)}{\log (1/\rho)}
   =\frac{2}{\eta\mu}\log\!\frac{1}{\varepsilon},
\]
so the overall computational complexity to reach $\varepsilon$‑accuracy is
\[
\mathcal{O}\!\Bigl(\bigl(n+m\bigr)\frac{1}{\eta\mu}\log\frac{1}{\varepsilon}\Bigr).
\]
Choosing $m=2n$ and $\eta=\frac{1}{3L}$ gives the classic SVRG complexity
\[
\mathcal{O}\!\Bigl(\bigl(n+2n\bigr)\frac{L}{\mu}\log\frac{1}{\varepsilon}\Bigr)
= \mathcal{O}\!\Bigl(\frac{L}{\mu}n\log\frac{1}{\varepsilon}\Bigr).
\]

\section*{Special Cases}
\begin{itemize}
\item \textbf{Exact gradients ($m=0$).} The algorithm reduces to full‑gradient descent with rate $(1-\eta\mu)$.
\item \textbf{Large inner loop ($m\to\infty$).} The term $\bigl(1-\eta\mu/2\bigr)^{m}$ vanishes, and the recursion (24) still holds, i.e. the method behaves like a deterministic gradient method after a single snapshot.
\item \textbf{Non‑strongly convex case ($\mu=0$).} The proof collapses; SVRG is known to achieve a sub‑linear $O(1/T)$ rate under mere convexity (see \cite{johnson2013accelerating}).
\end{itemize}

\begin{thebibliography}{9}
\bibitem{johnson2013accelerating}
R.~Johnson and T.~Zhang, ``Accelerating stochastic gradient descent using
  predictive variance reduction,'' \emph{Advances in Neural Information
  Processing Systems}, vol.~26, 2013.
\end{thebibliography}

\end{document}